\documentclass{article} 
\usepackage[T1]{fontenc}
\usepackage[utf8]{inputenc}
\usepackage{graphicx}       
\usepackage{geometry}       
\usepackage{listings}  
\usepackage{hyperref}       
\usepackage{titling}        
\usepackage{setspace}  % paquete para interlineado

\setstretch{1.5} 
\hypersetup{
    colorlinks= true,
    linkcolor= blue,
    urlcolor=blue
}
\geometry{margin=1.5cm}

\title{Análisis y especificaciones de requisitos}
\author{Iván Alfaya Pérez}
\date{\today}
\newcommand{\subtitle}{Proyecto: Simulación Cafetería JAVA/JavaFX}
\posttitle{\par\vspace{0.5cm}\large \subtitle\par\vskip 0.5cm}

\begin{document}
\maketitle

\clearpage 
\tableofcontents 
\clearpage 

%---------------------------------------------------------
\section{Introducción}
El proyecto tiene como objetivo simular el funcionamiento de una cafetería, con sus camareros y clientes.
Se utiliza JavaFX para la interfaz gráfica y \texttt{Threads} de Java para simular la concurrencia entre clientes y camareros.  
El sistema muestra en tiempo real el estado de los clientes y la preparación de sus pedidos, considerando la paciencia de cada cliente.

%---------------------------------------------------------
\section{Requisitos funcionales}
\begin{itemize}
    \item RF1: El sistema permite la creación de clientes con tiempos de espera (paciencia) definidos.
    \item RF2: Cada camarero atiende a varios clientes de forma concurrente.
    \item RF3: Si el tiempo de preparación del café supera la paciencia del cliente, este se retira sin ser atendido.
    \item RF4: La interfaz gráfica muestra el estado de cada cliente y camarero en tiempo real.
    \item RF5: Al finalizar la simulación, se muestra un resumen de los clientes atendidos y los que se marcharon.
\end{itemize}

%---------------------------------------------------------
\section{Requisitos no funcionales}
\begin{itemize}
    \item RNF1: La aplicación se ejecuta en tiempo real utilizando hilos de Java.
    \item RNF2: El código es modular y orientado a objetos.
    \item RNF3: La interfaz gráfica es intuitiva y clara.
    \item RNF4: El sistema es escalable para soportar más clientes y camareros.
\end{itemize}

%---------------------------------------------------------
\section{Casos de uso}
\begin{itemize}
    \item \textbf{CU1 - Llegada de Cliente:} El cliente llega a la cafetería y espera ser atendido.
    \item \textbf{CU2 - Atención de Cliente:} Los camareros preparan los cafés de manera concurrente.
    \item \textbf{CU3 - Cliente Impaciente:} Si el tiempo de espera es demasiado largo, el cliente se marcha.
    \item \textbf{CU4 - Finalización del Servicio:} Al terminar, se muestra un resumen de los clientes atendidos.
\end{itemize}



\newpage
\section{Estructura general del sistema}
El sistema se compone de varias clases principales: \texttt{HelloApplication}, \texttt{HelloController}, \texttt{Cliente} y \texttt{Camarero}.  
Se utiliza JavaFX para la interfaz gráfica y Threads para la simulación concurrente.  

\begin{table}[h!]
\centering
\begin{tabular}{|c|p{10cm}|}
\hline
\textbf{Clase / Componente} & \textbf{Descripción} \\
\hline
HelloApplication & Clase principal que inicia la aplicación JavaFX y carga la vista FXML. \\
\hline
HelloController & Controlador de la interfaz. Maneja eventos de botones y actualiza los TextArea con el estado de los camareros y clientes. \\
\hline
Cliente & Representa a un cliente. Cada cliente es un \texttt{Thread} que espera su café y puede irse si se excede su tiempo de paciencia. \\
\hline
Camarero & Representa un camarero. Cada camarero es un \texttt{Thread} que prepara cafés para los clientes asignados y actualiza la interfaz mediante \texttt{Platform.runLater()}. \\
\hline
Vista FXML & Interfaz visual con dos paneles para camareros, TextArea para mensajes y botón para iniciar la simulación. \\
\hline
\end{tabular}
\caption{Clases y componentes principales del proyecto JavaFX.}
\label{tab:clasesJavaFX}
\end{table}





\section{Diagrama conceptual simplificado}

\begin{verbatim}
HelloApplication
 └── Inicia la interfaz JavaFX y carga HelloController

HelloController
 ├── Maneja eventos (botón "Empezar")
 ├── Crea clientes y camareros
 └── Actualiza TextArea y Label con clientes atendidos

Cliente (Thread)
 ├── nombre
 ├── tiempoEspera
 ├── atendido
 └── run(): espera y puede ser interrumpido si se atiende

Camarero (Thread)
 ├── nombre
 ├── listaClientes
 ├── TextArea para mensajes
 ├── atendidos (contador)
 └── run(): prepara cafés y actualiza interfaz
\end{verbatim}

%---------------------------------------------------------
\section{Interfaz JavaFX}
La interfaz consiste en:

\begin{itemize}
    \item Un título principal “SIMULACIÓN CAFETERÍA”.
    \item Dos paneles para cada camarero, con un \texttt{TextArea} que muestra los mensajes de preparación y atención de clientes.
    \item Un Label que indica el número de clientes atendidos.
    \item Un botón “Empezar” para iniciar la simulación.
\end{itemize}

La interfaz permite ver en tiempo real el flujo de clientes y la actividad de los camareros, mostrando cuándo los clientes reciben su café o se retiran por impaciencia.

%---------------------------------------------------------
\end{document}
